% GENERATED FILE. Edit canonical lessons, then run npm run generate:books.
% canonical-source-hash: fnv1a64:6b553b39e59ba6a7
% canonical-lessons: ML-C563-kathapatram, ML-C563-ramgam, ML-C563-vedi, ML-C563-kani, ML-C563-kayyati

\chapter{On Stage}
\label{ch:ml-on-stage}
\hlchaptermodality{563}

\begin{chapteropening}
\textbf{By the end of this chapter:} \emph{I can say a story character, a stage scene, a stage, an onlooker and applause.}

Complete the last lesson of chapter 563: I can say a story character, a stage scene, a stage, an onlooker and applause.
\end{chapteropening}

\section[\texorpdfstring{\ml{കഥാപാത്രം} (kathāpātraṁ)}{കഥാപാത്രം (kathāpātraṁ)}]{\ml{കഥാപാത്രം} (kathāpātraṁ) — a story character}

\label{lesson:ML-C563-kathapatram}



\begin{quote}
Before the new one: say the Malayalam for a screenplay, then the Malayalam for a novel.
\end{quote}

\subsection*{You'll want to know: \ml{കഥാപാത്രം}}
\textbf{\ml{കഥാപാത്രം}} — \emph{kathāpātraṁ} — \textquotedblleft{}a story character\textquotedblright{}.

\begin{quote}
\ml{നോവലിലെ} \ml{പ്രധാന} \ml{കഥാപാത്രം} \ml{ഒരു} \ml{കർഷകനാണ്}.
\end{quote}

\begin{quote}
\emph{nōvalile pradhāna kathāpātraṁ oru karṣakanāṇŭ.}
\end{quote}

\begin{quote}
The main character of the novel is a farmer.
\end{quote}

\ml{കഥ} a story + \ml{പാത്രം} a vessel: the vessel that carries the story.

\begin{cousinweb}
\textbf{\ml{കഥാപാത്രം}} (\emph{kathāpātraṁ}) is a learned compound built from Sanskrit \emph{kathā}, \textquotedblleft{}a story\textquotedblright{} (CDIAL 2704), and \emph{pātra}, \textquotedblleft{}a vessel, a dish\textquotedblright{} (CDIAL 8055).
\end{cousinweb}

\begin{grammarlens}[title={What you've built}]
One more word for describing what you lived and felt.
\end{grammarlens}

\subsection*{Your turn}
\begin{itemize}
\raggedright
  \item \emph{Say it:} \emph{kathāpātraṁ}
  \item \emph{Say it:} \emph{kathāpātraṁ}, once more
  \item \emph{Say it:} \emph{nōvalile pradhāna kathāpātraṁ oru karṣakanāṇŭ.} — The main character of the novel is a farmer.
  \item \emph{Recall:} say the Malayalam for a screenplay, then the Malayalam for a novel, then say \emph{kathāpātraṁ} again
\end{itemize}

\begin{tcolorbox}[breakable,colback=teal!4,colframe=teal!35!black,title={Before you move on}]
What does \ml{കഥാപാത്രം} mean? (\textquotedblleft{}A story character\textquotedblright{}.) Say it once more.
\end{tcolorbox}
\section[\texorpdfstring{\ml{രംഗം} (raṁgaṁ)}{രംഗം (raṁgaṁ)}]{\ml{രംഗം} (raṁgaṁ) — a stage scene; a scene}

\label{lesson:ML-C563-ramgam}



\begin{quote}
Before the new one: say the Malayalam for a novel, then the Malayalam for a story character.
\end{quote}

\subsection*{You'll want to know: \ml{രംഗം}}
\textbf{\ml{രംഗം}} — \emph{raṁgaṁ} — \textquotedblleft{}a stage scene; a scene\textquotedblright{}.

\begin{quote}
\ml{സിനിമയിലെ} \ml{അവസാന} \ml{രംഗം} \ml{എല്ലാവരെയും} \ml{കരയിച്ചു}.
\end{quote}

\begin{quote}
\emph{sinimayile avasāna raṁgaṁ ellāvareyuṁ karayiccu.}
\end{quote}

\begin{quote}
The last scene of the film made everyone cry.
\end{quote}

Also a field of work: \ml{കലാരംഗം}, the world of the arts.

\begin{cousinweb}
\textbf{\ml{രംഗം}} (\emph{raṁgaṁ}) is borrowed from Sanskrit \emph{raṅga}, \textquotedblleft{}dancing, a place of public amusement, a stage\textquotedblright{} (CDIAL 10561).
\end{cousinweb}

\begin{grammarlens}[title={What you've built}]
One more word for describing what you lived and felt.
\end{grammarlens}

\subsection*{Your turn}
\begin{itemize}
\raggedright
  \item \emph{Say it:} \emph{raṁgaṁ}
  \item \emph{Say it:} \emph{raṁgaṁ}, once more
  \item \emph{Say it:} \emph{sinimayile avasāna raṁgaṁ ellāvareyuṁ karayiccu.} — The last scene of the film made everyone cry.
  \item \emph{Recall:} say the Malayalam for a novel, then the Malayalam for a story character, then say \emph{raṁgaṁ} again
\end{itemize}

\begin{tcolorbox}[breakable,colback=teal!4,colframe=teal!35!black,title={Before you move on}]
What does \ml{രംഗം} mean? (\textquotedblleft{}A stage scene; a scene\textquotedblright{}.) Say it once more.
\end{tcolorbox}
\section[\texorpdfstring{\ml{വേദി} (vēdi)}{വേദി (vēdi)}]{\ml{വേദി} (vēdi) — a stage; a venue}

\label{lesson:ML-C563-vedi}



\begin{quote}
Before the new one: say the Malayalam for a story character, then the Malayalam for a stage scene.
\end{quote}

\subsection*{You'll want to know: \ml{വേദി}}
\textbf{\ml{വേദി}} — \emph{vēdi} — \textquotedblleft{}a stage; a venue\textquotedblright{}.

\begin{quote}
\ml{ഗായകൻ} \ml{വേദിയിൽ} \ml{കയറിയപ്പോൾ} \ml{എല്ലാവരും} \ml{എഴുന്നേറ്റു}.
\end{quote}

\begin{quote}
\emph{gāyakan vēdiyil kayaṟiyappōḷ ellāvaruṁ eḻunnēṟṟu.}
\end{quote}

\begin{quote}
When the singer came on stage, everyone stood up.
\end{quote}

Any stage or meeting place: \ml{സമ്മേളന} \ml{വേദി}, a conference venue.

\begin{cousinweb}
\textbf{\ml{വേദി}} (\emph{vēdi}) is borrowed from Sanskrit \emph{vedi}, the raised ground prepared for a rite (Monier-Williams); a raised place for a rite became a raised place for a show.
\end{cousinweb}

\begin{grammarlens}[title={What you've built}]
One more word for describing what you lived and felt.
\end{grammarlens}

\subsection*{Your turn}
\begin{itemize}
\raggedright
  \item \emph{Say it:} \emph{vēdi}
  \item \emph{Say it:} \emph{vēdi}, once more
  \item \emph{Say it:} \emph{gāyakan vēdiyil kayaṟiyappōḷ ellāvaruṁ eḻunnēṟṟu.} — When the singer came on stage, everyone stood up.
  \item \emph{Recall:} say the Malayalam for a story character, then the Malayalam for a stage scene, then say \emph{vēdi} again
\end{itemize}

\begin{tcolorbox}[breakable,colback=teal!4,colframe=teal!35!black,title={Before you move on}]
What does \ml{വേദി} mean? (\textquotedblleft{}A stage; a venue\textquotedblright{}.) Say it once more.
\end{tcolorbox}
\section[\texorpdfstring{\ml{കാണി} (kāṇi)}{കാണി (kāṇi)}]{\ml{കാണി} (kāṇi) — an onlooker; a spectator}

\label{lesson:ML-C563-kani}



\begin{quote}
Before the new one: say the Malayalam for a stage scene, then the Malayalam for a stage.
\end{quote}

\subsection*{You'll want to know: \ml{കാണി}}
\textbf{\ml{കാണി}} — \emph{kāṇi} — \textquotedblleft{}an onlooker; a spectator\textquotedblright{}.

\begin{quote}
\ml{വഴിയിൽ} \ml{നടന്ന} \ml{വഴക്ക്} \ml{കാണാൻ} \ml{ഒരുപാട്} \ml{കാണികൾ} \ml{കൂടി}.
\end{quote}

\begin{quote}
\emph{vaḻiyil naṭanna vaḻakkŭ kāṇān orupāṭŭ kāṇikaḷ kūṭi.}
\end{quote}

\begin{quote}
Lots of onlookers gathered to watch the quarrel in the street.
\end{quote}

Plural: \ml{കാണികൾ}, \emph{kāṇikaḷ}.

Plural: \ml{കാണികൾ}, the spectators at a match or a show.

\begin{cousinweb}
\textbf{\ml{കാണി}} (\emph{kāṇi}) is native Dravidian: DEDR 1443 lists it with \ml{കാണുക}, \textquotedblleft{}to see\textquotedblright{}, as a visitor or spectator (DEDR 1443).
\end{cousinweb}

\begin{grammarlens}[title={What you've built}]
One more word for describing what you lived and felt.
\end{grammarlens}

\subsection*{Your turn}
\begin{itemize}
\raggedright
  \item \emph{Say it:} \emph{kāṇi}
  \item \emph{Say it:} \emph{kāṇi}, once more
  \item \emph{Say it:} \emph{vaḻiyil naṭanna vaḻakkŭ kāṇān orupāṭŭ kāṇikaḷ kūṭi.} — Lots of onlookers gathered to watch the quarrel in the street.
  \item \emph{Recall:} say the Malayalam for a stage scene, then the Malayalam for a stage, then say \emph{kāṇi} again
\end{itemize}

\begin{tcolorbox}[breakable,colback=teal!4,colframe=teal!35!black,title={Before you move on}]
What does \ml{കാണി} mean? (\textquotedblleft{}An onlooker; a spectator\textquotedblright{}.) Say it once more.
\end{tcolorbox}
\section[\texorpdfstring{\ml{കയ്യടി} (kayyaṭi)}{കയ്യടി (kayyaṭi)}]{\ml{കയ്യടി} (kayyaṭi) — applause}

\label{lesson:ML-C563-kayyati}



\begin{quote}
Before the new one: say the Malayalam for a stage, then the Malayalam for an onlooker.
\end{quote}

\subsection*{You'll want to know: \ml{കയ്യടി}}
\textbf{\ml{കയ്യടി}} — \emph{kayyaṭi} — \textquotedblleft{}applause\textquotedblright{}.

\begin{quote}
\ml{പാട്ട്} \ml{കഴിഞ്ഞപ്പോൾ} \ml{വലിയ} \ml{കയ്യടി} \ml{ഉണ്ടായി}.
\end{quote}

\begin{quote}
\emph{pāṭṭŭ kaḻiññappōḷ valiya kayyaṭi uṇṭāyi.}
\end{quote}

\begin{quote}
There was loud applause when the song ended.
\end{quote}

\ml{കൈ} hand + \ml{അടി} a stroke; \ml{കയ്യടിക്കുക} is to applaud.

\begin{cousinweb}
\textbf{\ml{കയ്യടി}} (\emph{kayyaṭi}) is native Dravidian, a compound of \ml{കൈ}, \textquotedblleft{}hand\textquotedblright{} (DEDR 2023), and \ml{അടി}, \textquotedblleft{}a blow, a stroke\textquotedblright{} (DEDR 77).
\end{cousinweb}

\begin{grammarlens}[title={What you've built}]
One more word for describing what you lived and felt.
\end{grammarlens}

\subsection*{Your turn}
\begin{itemize}
\raggedright
  \item \emph{Say it:} \emph{kayyaṭi}
  \item \emph{Say it:} \emph{kayyaṭi}, once more
  \item \emph{Say it:} \emph{pāṭṭŭ kaḻiññappōḷ valiya kayyaṭi uṇṭāyi.} — There was loud applause when the song ended.
  \item \emph{Recall:} say the Malayalam for a stage, then the Malayalam for an onlooker, then say \emph{kayyaṭi} again
\end{itemize}

\begin{tcolorbox}[breakable,colback=teal!4,colframe=teal!35!black,title={Before you move on}]
What does \ml{കയ്യടി} mean? (\textquotedblleft{}Applause\textquotedblright{}.) Say it once more.
\end{tcolorbox}
